\documentclass[../../../ps_q.tex]{subfiles}

\begin{document}
    \begin{enumerate}
        \item 図のように，電池$\mathrm{E}_1$（起電力$3V_0$），$\mathrm{E}_2$（起電力$V_0$），%
            電気抵抗$\mathrm{R}_1$（抵抗値$2R$），%
            $\mathrm{R}_2$（抵抗値$2R$），$\mathrm{R}_3$（抵抗値$R$）からなる回路がある．%
            AB間を流れる電流の大きさ$I$を求めよ．\\
            \begin{minipage}{\linewidth}
                \centering
                    \includegraphics{\subfix{fig_em_cb_08/fig_em_cb_08_01_q.pdf}}
            \end{minipage}
            \vspace{1pt}
        \item 図のように，電池$\mathrm{E}$（起電力$V_0$），電気抵抗$\mathrm{R}_1$（抵抗値$R$），%
            $\mathrm{R}_2$（抵抗値$R$），$\mathrm{R}_3$（抵抗値$R$），%
            $\mathrm{R}_4$（抵抗値$3R$），$\mathrm{R}_5$（抵抗値$R$）からなる回路がある．%
            AB間を流れる電流の大きさ$I$を求めよ．\\        
            \begin{minipage}{\linewidth}
                \centering
                    \includegraphics{\subfix{fig_em_cb_08/fig_em_cb_08_02_q.pdf}}
            \end{minipage}
    \end{enumerate}
\end{document}




